problem:  
Let \(M^6\) be a fixed compact, simply connected, smooth manifold that admits at least one strict nearly Kähler structure \((J, g)\).  
Define the *moduli space* \(\mathcal{N}(M)\) as the set of strict nearly Kähler SU(3)-structures on \(M\) modulo the action of the diffeomorphism group isotopic to the identity.  

It is known (from work of Butruille and others) that for homogeneous nearly Kähler structures, infinitesimal deformations correspond to certain Laplacian eigenforms, and that these examples are analytically “rigid.” However, the discovery of inhomogeneous nearly Kähler structures (e.g., by Foscolo–Haskins on \(S^6\) and \(S^3\times S^3\)) suggests a richer deformation space.  

**Problem:**  
For a given compact simply connected \(6\)-manifold \(M\) that admits a strict nearly Kähler structure, determine whether \(\mathcal{N}(M)\) is discrete or contains nontrivial (possibly smooth) components. In particular:  
- For \(M = S^6\) and \(M = S^3\times S^3\), does there exist a continuous 1-parameter family of pairwise non-isometric strict nearly Kähler structures?  
- Is it possible that \(\mathcal{N}(M)\) possesses distinct connected components corresponding to inequivalent torsion classes or topological data (e.g., cohomological invariants of the SU(3)-structure)?  

Provide possible geometric or analytical mechanisms (Laplace-type spectral conditions, stability or obstruction theories, including those derived from the Hitchin functional or G_2-cone deformation theory) that could govern the existence or discreteness of such moduli.

Why is it a "good" problem:  
This question is deeply and precisely formulated within the current frontier of differential geometry. It fixes the manifold and asks about the structure of its nearly Kähler moduli space, a problem whose answer is genuinely unknown but tightly connected with open analytical and topological obstructions. It probes the rigidity versus flexibility of nearly Kähler SU(3)-structures, integrating ideas from nonlinear analysis (PDE systems defining torsion classes), topology (diffeomorphism invariants and characteristic classes), and even higher-dimensional geometry (via the relation to G_2-cones). The problem is simple to state—“Are strict nearly Kähler structures isolated or deformable?”—yet profoundly deep, serving as a bridge between geometry, analysis, and topology, and capable of guiding future classification results.